% Part: first-order-logic
% Chapter: completeness
% Section: henkin-expansion

\documentclass[../../../include/open-logic-section]{subfiles}

\begin{document}

\olfileid[fr]{fol}{com}{hen}

\olsection{Expansion de Henkin}

\begin{explain}
Une difficulté de la démonstration de complétude est que le modèle
construit à partir d'un ensemble cohérent et complet~$\Gamma$ doit
rendre vraies toutes les !!{formula}s quantifiées de~$\Gamma$.
Pour le garantir, nous recourons à une méthode due à Leon Henkin.
Elle consiste essentiellement à enrichir le langage d'une infinité
de nouvelles !!{constant}s et à ajouter, pour chaque !!{formula}
$!A(x)$ dont les variables libres appartiennent à~$\{x\}$, une
!!{formula} de la forme
\iftag{prvEx}
      {$\lexists[x][!A(x)] \lif !A(c)$}
      {$\lnot\lforall[x][!A(x)] \lif \lnot !A(c)$},
où $c$ est l'une des nouvelles !!{constant}s. Dans la !!{structure}
que nous construirons,
\iftag{prvEx}
{tout !!{sentence} existentiel vrai aura un témoin}
{tout !!{sentence} universel faux aura un contre-exemple}
parmi ces constantes.
\end{explain}

\begin{prop}
\ollabel{prop:lang-exp}
Si $\Gamma$ est cohérent dans $\Lang L$ et si $\Lang L'$ est obtenu
en ajoutant à $\Lang L$ un ensemble infini dénombrable de nouvelles
!!{constant}s $\Obj d_0$, $\Obj d_1$, \dots, alors $\Gamma$ reste
cohérent dans~$\Lang L'$.
\end{prop}

\begin{proof}
Pour comparer les deux langages, on dispose, dans les étapes
intermédiaires des dérivations, d'une réserve infinie de constantes
auxiliaires distinctes des constantes du langage des hypothèses
et de la conclusion. Supposons qu'une dérivation de l'absurde
à partir de~$\Gamma$ existe dans~$\Lang L'$. Elle n'utilise qu'un
nombre fini de nouvelles constantes. Renommons-les uniformément
par des constantes auxiliaires distinctes, choisies en dehors
de celles qui figurent déjà dans la dérivation. Ce renommage
préserve les instances des règles et leurs conditions de fraîcheur.
Les hypothèses et l'absurde sont inchangés. Nous obtenons une
dérivation de l'absurde à partir de~$\Gamma$ dans~$\Lang L$,
ce qui contredit sa cohérence.
\end{proof}

\begin{defn}[Ensemble saturé]
Un ensemble~$\Gamma$ de !!{formula}s d'un langage~$\Lang{L}$ est
\emph{saturé} si et seulement si, pour chaque !!{formula}
$!A(x)\in\Frm[L]$ dont les !!{variable}s libres appartiennent
à~$\{x\}$, il existe !!a{constant}~$c\in\Lang{L}$ telle que
\iftag{prvEx}
      {$\lexists[x][!A(x)] \lif !A(c) \in \Gamma$}
      {$\lnot\lforall[x][!A(x)] \lif \lnot !A(c) \in \Gamma$}.
\end{defn}

Nous utiliserons la définition suivante pour démontrer le résultat
suivant.

\begin{defn}
\ollabel{defn:henkin-exp}
Soit $\Lang L'$ le langage de la \olref{prop:lang-exp}. Fixons une
énumération $!A_0(x_0)$, $!A_1(x_1)$, \dots des couples formés
d'une !!{formula}~$!A_i(x_i)$ de~$\Lang L'$ et d'une
!!{variable}~$x_i$ telle que l'ensemble des variables libres de
la formule soit inclus dans~$\{x_i\}$. Comme le langage est au plus dénombrable,
ces couples sont dénombrables. Nous définissons les
!!{sentence}s~$!D_n$ par récurrence sur~$n$.

Soit $c_0$ la première des !!{constant}s $\Obj d_i$ ajoutées à
$\Lang{L}$ qui ne figure pas dans~$!A_0(x_0)$. Supposons
$!D_0,\dots,!D_{n-1}$ déjà définies, et choisissons pour~$c_n$
la première des nouvelles !!{constant}s~$\Obj d_i$ qui ne figure
ni dans $!D_0,\dots,!D_{n-1}$, ni dans~$!A_n(x_n)$.

Définissons alors la !!{formula}~$!D_n$ par
\iftag{prvEx}
{$\lexists[x_{n}][!A_{n}(x_{n})] \lif
  !A_{n}(c_{n})$}{$\lnot\lforall[x_{n}][!A_{n}(x_{n})] \lif \lnot
  !A_{n}(c_{n})$}.
\end{defn}

\begin{lem}
\ollabel{lem:henkin}
Tout ensemble cohérent~$\Gamma$ peut s'étendre en un ensemble
cohérent et saturé~$\Gamma'$.
\end{lem}

\begin{proof}
Partons d'un ensemble cohérent d'!!{sentence}s~$\Gamma$ dans un
langage~$\Lang{L}$ et enrichissons ce langage d'un ensemble infini
dénombrable de nouvelles !!{constant}s pour former~$\Lang{L'}$.
Par la \olref{prop:lang-exp}, $\Gamma$ reste cohérent dans ce langage.
Soient ensuite les $!D_i$ de la \olref{defn:henkin-exp}. Posons
\begin{align*}
\Gamma_0 & = \Gamma \\
\Gamma_{n+1} & = \Gamma_n \cup \{!D_n \}
\end{align*}
autrement dit $\Gamma_{n+1}=\Gamma\cup\{!D_0,\dots,!D_n\}$,
puis $\Gamma'=\bigcup_n\Gamma_n$. Par construction, $\Gamma'$
est saturé.

Si $\Gamma'$ était incohérent, un certain $\Gamma_n$ le serait aussi
(exercice : expliquez pourquoi). Il suffit donc de montrer par
récurrence sur~$n$ que tous les $\Gamma_n$ sont cohérents.

Le cas de base est l'hypothèse même du lemme :
$\Gamma_0=\Gamma$ est cohérent. Pour l'étape de récurrence,
supposons $\Gamma_n$ cohérent, mais
$\Gamma_{n+1}=\Gamma_n\cup\{!D_n\}$ incohérent. Rappelons que
$!D_n$ vaut
\iftag{prvEx}
{$\lexists[x_{n}][!A_{n}(x_n)] \lif !A_{n}(c_{n})$}
{$\lnot\lforall[x_{n}][!A_{n}(x_n)] \lif \lnot !A_{n}(c_{n})$},
où $!A_n(x_n)$ est une !!{formula} de~$\Lang{L'}$ dont les
!!{variable}s libres appartiennent à~$\{x_n\}$. Par le choix
de~$c_n$ dans la
\olref{defn:henkin-exp}, cette constante ne figure ni dans
$!A_n(x_n)$, ni dans~$\Gamma_n$.

Si $\Gamma_n\cup\{!D_n\}$ est incohérent, alors
$\Gamma_n\Proves\lnot !D_n$. La négation de l'implication $!D_n$
donne son antécédent et la négation de son conséquent%
\iftag{prvEx}{.}{, puis la double négation du conséquent s'élimine.}
\iftag{prvEx}{
\[
\Gamma_n \Proves \lexists[x_n][!A_n(x_n)]
\qquad
\Gamma_n \Proves \lnot !A_n(c_n)
\]}{
\[
\Gamma_n \Proves \lnot\lforall[x_n][!A_n(x_n)]
\qquad
\Gamma_n \Proves !A_n(c_n)
\]}
Puisque $c_n$ ne figure ni dans $\Gamma_n$, ni dans~$!A_n(x_n)$,
les \tagrefs{prfAX/{fol:axd:qpr:thm:strong-generalization},prfSC/{fol:seq:qpr:thm:strong-generalization},prfND/{fol:ntd:qpr:thm:strong-generalization},prfTab/{fol:tab:qpr:thm:strong-generalization}}
s'appliquent. De
\iftag{prvEx}
{$\Gamma_n \Proves \lnot !A_n(c_n)$}
{$\Gamma_n \Proves !A_n(c_n)$},
nous déduisons
\iftag{prvEx}
{$\Gamma_n \Proves \lforall[x_n][\lnot !A_n(x_n)]$}
{$\Gamma_n \Proves \lforall[x_n][!A_n(x_n)]$}.
Ainsi, on aurait simultanément
\iftag{prvEx}
{$\Gamma_n \Proves \lexists[x_n][!A_n(x_n)]$ et
$\Gamma_n \Proves \lforall[x_n][\lnot !A_n(x_n)]$}
{$\Gamma_n \Proves \lnot\lforall[x_n][!A_n(x_n)]$ et
$\Gamma_n \Proves \lforall[x_n][!A_n(x_n)]$},
ce qui rendrait $\Gamma_n$ incohérent.
\iftag{prvEx}
{Notons que
\iftag{prvAll}
{$\lforall[x_n][\lnot !A_n(x_n)] \Proves
  \lnot\lexists[x_n][!A_n(x_n)]$}
{$\lforall[x_n][\lnot !A_n]$ est défini par
  $\lnot\lexists[x_n][\lnot\lnot !A_n(x_n)]$ et que
  $\lnot\lexists[x_n][\lnot\lnot !A_n(x_n)] \Proves
  \lnot\lexists[x_n][!A_n(x_n)]$}.}{}
Contradiction : $\Gamma_n$ était supposé cohérent. Donc
$\Gamma_n\cup\{!D_n\}$ est cohérent.
\end{proof}

\begin{explain}
Dans un ensemble cohérent, !!{complete} et saturé, l'appartenance
d'un !!{sentence} quantifié dépend de celle de ses instances closes.
\iftag{prvAll}{Un !!{sentence} universel appartient à l'ensemble si
et seulement si toutes ses instances closes lui appartiennent.}{}
\iftag{prvEx}{Un !!{sentence} existentiel appartient à l'ensemble si
et seulement si au moins une instance close lui appartient.}{}
Nous en déduirons que la !!{structure} construite à partir d'un tel
ensemble rend vrais tous ses !!{sentence}s quantifiés.
\end{explain}

\begin{prop}\ollabel{prop:saturated-instances}
Supposons $\Gamma$ complet, cohérent et saturé.
\begin{tagenumerate}{prvEx,prvAll}
\tagitem{prvEx}{$\lexists[x][!A(x)] \in \Gamma$ si et seulement si
  $!A(t) \in \Gamma$ pour au moins un terme clos~$t$.}{}
\tagitem{prvAll}{$\lforall[x][!A(x)] \in \Gamma$ si et seulement si
  $!A(t) \in \Gamma$ pour tout terme clos~$t$.}{}
\end{tagenumerate}
\end{prop}

\begin{proof}
  \begin{tagenumerate}{prvEx,prvAll}
  \tagitem{prvEx}{%
    \iftag{probEx}{Exercice.}{Supposons d'abord
      $\lexists[x][!A(x)]\in\Gamma$. Comme $\Gamma$ est saturé,
      $(\lexists[x][!A(x)]\lif !A(c))\in\Gamma$ pour une
      !!{constant}~$c$. D'après les
      \tagrefs{prfAX/{fol:axd:ppr:prop:provability-lif},prfSC/{fol:seq:ppr:prop:provability-lif},prfND/{fol:ntd:ppr:prop:provability-lif},prfTab/{fol:tab:ppr:prop:provability-lif}}, point~(1),
      et la \olref[ccs]{prop:ccs}\olref[ccs]{prop:ccs-prov-in},
      $!A(c)\in\Gamma$.

    Pour la réciproque, la saturation n'est pas nécessaire.
    Supposons $!A(t)\in\Gamma$. Alors
    $\Gamma\Proves\lexists[x][!A(x)]$ d'après les
      \tagrefs{prfAX/{fol:axd:qpr:prop:provability-quantifiers},prfSC/{fol:seq:qpr:prop:provability-quantifiers},prfND/{fol:ntd:qpr:prop:provability-quantifiers},prfTab/{fol:tab:qpr:prop:provability-quantifiers}}, point~(1).
    Par la \olref[ccs]{prop:ccs}\olref[ccs]{prop:ccs-prov-in},
    $\lexists[x][!A(x)]\in\Gamma$.}}{}
  \tagitem{prvAll}{%
    \iftag{probAll}{Exercice.}{Supposons que $!A(t)\in\Gamma$ pour
      tout terme clos~$t$. Par l'absurde, supposons
      $\lforall[x][!A(x)]\notin\Gamma$. Puisque $\Gamma$ est
      complet, $\lnot\lforall[x][!A(x)]\in\Gamma$. Par saturation,
      \iftag{prvEx}{$(\lexists[x][\lnot !A(x)]\lif\lnot !A(c))
        \in\Gamma$}{$(\lnot\lforall[x][!A(x)]\lif\lnot !A(c))
        \in\Gamma$} pour une !!{constant}~$c$. Par hypothèse,
      $!A(c)\in\Gamma$, car $c$ est un terme clos.
      \iftag{prvEx}{Par dualité classique, $\lnot\lforall[x][!A(x)]$
      donne $\lexists[x][\lnot !A(x)]$.}{}
      Ces énoncés rendraient $\Gamma$ incohérent. (Exercice : donnez
      la !!{derivation} montrant que
      \[
      \lnot \lforall[x][!A(x)],
      \iftag{prvEx}{\lexists[x][\lnot !A(x)] \lif \lnot
        !A(c)}{\lnot\lforall[x][!A(x)] \lif \lnot
        !A(c)}, !A(c)
      \]
      est incohérent.)

      Pour la réciproque, la saturation n'est pas nécessaire.
      Supposons $\lforall[x][!A(x)]\in\Gamma$. Alors
      $\Gamma\Proves !A(t)$ d'après les
      \tagrefs{prfAX/{fol:axd:qpr:prop:provability-quantifiers},prfSC/{fol:seq:qpr:prop:provability-quantifiers},prfND/{fol:ntd:qpr:prop:provability-quantifiers},prfTab/{fol:tab:qpr:prop:provability-quantifiers}},
      point~(2). Par la \olref[ccs]{prop:ccs}\olref[ccs]{prop:ccs-prov-in},
      $!A(t)\in\Gamma$.}}{}
  \end{tagenumerate}
\end{proof}

\end{document}
